% coords.tex -- the Araujo-Pardo--Lopez construction and the four-hat construction in projective
% coordinates of PG(2,q), with the degree lemma and the diameter theorem.  Drop-in section: needs
% amsmath, amssymb, amsthm (theorem, lemma, remark environments), enumitem, booktabs.
% Computer check: coords.py (builds the graphs only from this text), compare.py.
\providecommand{\PG}{\mathrm{PG}}
\providecommand{\F}{\mathbb{F}}
\providecommand{\cS}{\mathcal{S}}
\providecommand{\cA}{\mathcal{A}}
\newcommand{\lam}{\lambda}
\newcommand{\pnt}[3]{({#1}{:}{#2}{:}{#3})}
\newcommand{\lin}[3]{[{#1}{:}{#2}{:}{#3}]}

\section{The construction in coordinates}\label{sec:coords}

\subsection{Notation}\label{ssec:notation}
Let $q$ be a prime power and $\F = \mathrm{GF}(q)$, $\F^* = \F \setminus \{0\}$. A point of $\PG(2,q)$ is written $\pnt xyz$ and a line $\lin uvw$, with $(x,y,z), (u,v,w) \in \F^3 \setminus \{0\}$ up to nonzero scalars; the point $\pnt xyz$ is on the line $\lin uvw$ if and only if $ux + vy + wz = 0$. We fix
\[
\ell = \lin 100 \ (\text{the line } x = 0), \qquad A = \pnt 100, \qquad A \notin \ell,
\]
and name the points of $\ell$ and the lines through $A$:
\[
X_t = \pnt 0t1 \ (t \in \F), \quad P = \pnt 010; \qquad
\lam_t = \lin 01{-t} \ (t \in \F), \quad \lam_\infty = \lin 001 .
\]
Every point and every line has exactly one representative in the following lists.
\begin{center}
\begin{tabular}{@{}ll@{\qquad}ll@{}}
\toprule
points & number & lines & number\\
\midrule
$A = \pnt 100$ & $1$ & $\ell = \lin 100$ & $1$ \\
$X_t = \pnt 0t1$, $t \in \F$ & $q$ & $\lam_t = \lin 01{-t}$, $t \in \F$ & $q$ \\
$P = \pnt 010$ & $1$ & $\lam_\infty = \lin 001$ & $1$ \\
$\pnt xy1$, $x \in \F^*$, $y \in \F$ & $q(q-1)$ & $\lin u1w$, $u \in \F^*$, $w \in \F$ & $q(q-1)$\\
$\pnt x10$, $x \in \F^*$ & $q-1$ & $\lin 10w$, $w \in \F^*$ & $q-1$ \\
\bottomrule
\end{tabular}
\end{center}
The points not on $\ell$ are those with $x \ne 0$; the lines not through $A$ are those with $u \neq 0$. We use the following facts, each a one-line computation.
\begin{enumerate}[label=\textup{(F\arabic*)},leftmargin=3em]
\item\label{f:lam} $\lam_t \cap \ell = X_t$ and $\lam_\infty \cap \ell = P$. The points of $\lam_t$ are $A$ and $\pnt xt1$, $x \in \F$ (with $X_t$ at $x = 0$); the points of $\lam_\infty$ are $A$ and $\pnt x10$, $x \in \F$ (with $P$ at $x = 0$). In particular a point $\pnt xy1$ with $x \ne 0$ lies on $\lam_y$, and a point $\pnt x10$ lies on $\lam_\infty$.
\item\label{f:X} The lines through $X_t$ are $\ell$ and $\lin w1{-t}$, $w \in \F$ (with $\lam_t$ at $w = 0$). The lines through $P$ are $\ell$, $\lam_\infty$ and $\lin 10w$, $w \in \F^*$.
\item\label{f:join} If $u = \pnt xyz$ with $x \neq 0$ and $s \in \F$, then $u$ and $X_s$ lie on $\lin{(sz-y)/x}1{-s}$, which is not $\ell$.
\item\label{f:meet} A line $L = \lin uvw$ with $u \neq 0$ meets $\lam_s$ in $\pnt{-(vs+w)/u}s1$, which is not $A$; and $\lin u1w$ (with $u \ne 0$) meets $\ell$ in $X_{-w}$.
\end{enumerate}

\paragraph{A correlation.} Let $\delta$ map the point $\pnt xyz$ to the line $\lin xz{-y}$ and the line $\lin uvw$ to the point $\pnt uw{-v}$. Since $u x + w z + (-v)(-y) = ux + vy + wz$, $\delta$ maps incident pairs to incident pairs; it is a bijection from points to lines and from lines to points, and
\[
\delta(A) = \ell,\quad \delta(\ell) = A,\quad \delta(X_t) = \lam_t,\quad \delta(\lam_t) = X_t,\quad \delta(P) = \lam_\infty,\quad \delta(\lam_\infty) = P.
\]
Moreover $\pnt xyz$ lies on $\delta\pnt xyz$ if and only if $x^2 + zy - yz = x^2 = 0$, that is, if and only if the point lies on $\ell$. We use $\delta$ only for the final matching (Step~\ref{c:fill}) and in Remark~\ref{rem:duality}; the proofs below do not depend on it.

\paragraph{Dictionary.} In the affine coordinates of~\cite{AL}, $(x,y) = \pnt xy1$, $P_m = \pnt 1m0$, $P = \pnt 010$, $[m,b] = \lin m{-1}b$ (the line $y = mx+b$), $L_c = \lin 10{-c}$ (the line $x = c$) and $L = \lin 001$. Hence
\[
L_0 = \ell,\quad P_0 = A,\quad (0,s) = X_s,\quad [0,s] = \lam_s,\quad L = \lam_\infty,
\]
and the point $P$ is the same. In the coordinate-free description of the four-hat construction, $AX_s = \lam_s$, $AP = \lam_\infty$, and $\ell^* = \{P\} \cup \{X_t : t \notin \{0,1,c\}\}$.

\subsection{The construction of Araujo-Pardo and L\'opez}\label{ssec:al}
Let $q \ge 3$, fix $a \in \F \setminus \{0,1\}$ and put $S_{\mathrm{AL}} = \{0,1,a\}$. Add the points $Q, H_0, H_1, H_a$ and the lines $M, K_0, K_1, K_a$ (in~\cite{AL}: $Q_0$, $\widehat{(0,s)}$, $M_0$, $\widehat{[0,s]}$). The structure $\cA$ has these incidences:
\begin{enumerate}[label=\textup{(A\arabic*)},leftmargin=3em]
\item\label{a:plane} every incidence of $\PG(2,q)$ except $X_a \in \lam_a$;
\item\label{a:hat} $H_s$ ($s \in S_{\mathrm{AL}}$) lies on $\lin w1{-s}$ for every $w \in \F$;
\item\label{a:hatline} $K_s$ ($s \in S_{\mathrm{AL}}$) contains $\pnt xs1$ for every $x \in \F$;
\item\label{a:M} $M$ contains $P$, $X_t$ for $t \in \F \setminus \{0,1\}$, and $H_0, H_1, H_a$;
\item\label{a:Q} $Q$ lies on $\lam_\infty$, on $\lam_t$ for $t \in \F \setminus \{0,1\}$, and on $K_0, K_1, K_a$;
\item\label{a:links} $A \in K_a$, \ $H_a \in \ell$, \ $H_0 \in K_1$, \ $H_1 \in K_0$.
\end{enumerate}
By~\ref{f:X}, $H_s$ lies on the $q$ lines through $X_s$ other than $\ell$; by~\ref{f:lam}, $K_s$ contains the $q$ points of $\lam_s$ other than $A$. In the notation of~\cite{AL}, \ref{a:hat} is $\widehat{(0,s)} \sim [t,s]$ (as $[t,s] = \lin{-t}1{-s}$) and~\ref{a:hatline} is $\widehat{[0,s]} \sim (t,s)$.

\emph{Neighbour sets.} The new vertices of $\cA$ have these neighbours, each $q+2$ of them:
\begin{center}
\begin{tabular}{@{}ll@{\qquad}ll@{}}
$Q$: & $\lam_\infty$, $\lam_t$ ($t \in \F\setminus\{0,1\}$), $K_0, K_1, K_a$ & $M$: & $P$, $X_t$ ($t \in \F\setminus\{0,1\}$), $H_0, H_1, H_a$\\
$H_0$: & $\lin w10$ ($w \in \F$), $M$, $K_1$ & $K_0$: & $\pnt x01$ ($x \in \F$), $Q$, $H_1$\\
$H_1$: & $\lin w1{-1}$ ($w \in \F$), $M$, $K_0$ & $K_1$: & $\pnt x11$ ($x \in \F$), $Q$, $H_0$\\
$H_a$: & $\lin w1{-a}$ ($w \in \F$), $M$, $\ell$ & $K_a$: & $\pnt xa1$ ($x \in \F$), $Q$, $A$
\end{tabular}
\end{center}
The vertices of $\PG(2,q)$ whose neighbourhood changes are: $A$ ($+K_a$); $P$ and $X_t$, $t \notin \{0,1,a\}$ ($+M$); $X_a$ ($-\lam_a$, $+M$, $+K_a$); $X_0$, $X_1$ ($+K_0$, resp.\ $+K_1$); $\pnt xs1$, $x \neq 0$, $s \in S_{\mathrm{AL}}$ ($+K_s$); $\ell$ ($+H_a$); $\lam_\infty$ and $\lam_t$, $t \notin \{0,1,a\}$ ($+Q$); $\lam_a$ ($-X_a$, $+Q$, $+H_a$); $\lam_0$, $\lam_1$ ($+H_0$, resp.\ $+H_1$); $\lin u1{-s}$, $u \ne 0$, $s \in S_{\mathrm{AL}}$ ($+H_s$). Hence every vertex has degree $q+2$ except the points $\pnt xy1$ ($x \ne 0$, $y \notin S_{\mathrm{AL}}$) and $\pnt x10 = P_{1/x}$ ($x \ne 0$), and the lines $\lin x{-1}y = [x,y]$ ($x \ne 0$, $y \notin S_{\mathrm{AL}}$) and $\lin 10{-i} = L_i$ ($i \ne 0$), which have degree $q+1$; there are $(q-1)(q-2)$ of each. The matching of~\cite{AL} is
\[
(x,y) \sim [x,y] \ \ (x \neq 0,\ y \notin S_{\mathrm{AL}}), \qquad P_i \sim L_i,
\]
non-incident pairs because $(x,y) \in [x,y]$ iff $x^2 = 0$ and $P_i \notin L_i$ always. In~\cite{AL} the second family is indexed by $i \in \F \setminus \{0,1,a\}$; then $P_1, L_1, P_a, L_a$ keep degree $q+1$. With $i \in \F^*$ the graph is $(q+2)$-regular of order $2(q^2+q+5)$. The map $\delta$ of \S\ref{ssec:notation} gives another such matching, $p \sim \delta(p)$, as in Step~\ref{c:fill} below. (The maximum degree and the diameter do not depend on the matching.)

\emph{Diameter of $\cA$.} The proofs of Lemmas~\ref{lem:coords-points} and~\ref{lem:coords-lines} below apply to $\cA$ with these changes: in (P1) and (L1) only $s = a$ occurs; the pairs $\{H_0, X_1\}$, $\{H_1, X_0\}$, $\{Q, H_0\}$, $\{Q, H_1\}$ use $K_1$, $K_0$, $K_1$, $K_0$; the pairs $\{K_0, \lam_1\}$, $\{K_1, \lam_0\}$ use $H_1$, $H_0$; the pairs $\{M, K_0\}$, $\{M, K_1\}$ use $H_1$, $H_0$. All other cases are word for word the same with $T$ replaced by $\{0,1\}$. So $\cA$ also has diameter $3$, for every prime power $q \ge 3$.

\subsection{The four-hat construction}\label{ssec:four}
Let $q \ge 4$ and fix distinct $a, c \in \F \setminus \{0,1\}$ (this needs $q \ge 4$). Put
\[
S = \{0, 1, a, c\}, \qquad T = \{0, 1, c\}, \qquad \sigma\colon\ 0 \mapsto c \mapsto 1 \mapsto 0 \quad (\text{a $3$-cycle on } T).
\]
Add five points $Q, H_0, H_1, H_a, H_c$ and five lines $M, K_0, K_1, K_a, K_c$. The structure $\cS = \cS(q;a,c)$ has exactly the following incidences.
\begin{enumerate}[label=\textup{(C\arabic*)},leftmargin=3em]
\item\label{c:plane} Every incidence of $\PG(2,q)$ except the four incidences $X_s \in \lam_s$, $s \in S$; and the three new incidences $X_t \in \lam_{\sigma(t)}$, $t \in T$:
\[
X_0 \in \lam_c, \qquad X_c \in \lam_1, \qquad X_1 \in \lam_0 .
\]
\item\label{c:hat} For $s \in S$, $H_s$ lies on $\lin w1{-s}$ for every $w \in \F$.
\item\label{c:hatline} For $s \in S$, $K_s$ contains $\pnt xs1$ for every $x \in \F$.
\item\label{c:M} $M$ contains $P$, $X_t$ for every $t \in \F \setminus T$, and $H_0, H_1, H_a, H_c$.
\item\label{c:Q} $Q$ lies on $\lam_\infty$, on $\lam_t$ for every $t \in \F \setminus T$, and on $K_0, K_1, K_a, K_c$.
\item\label{c:links} $A \in K_a$, \ $H_a \in \ell$, and $H_t \in K_{\sigma(t)}$ for $t \in T$:
\[
H_0 \in K_c, \qquad H_c \in K_1, \qquad H_1 \in K_0 .
\]
\item\label{c:fill} \textup{(Optional filler.)} Join each point $p$ in $\{\pnt xy1 : x \neq 0,\ y \notin S\} \cup \{\pnt x10 : x \neq 0\}$ to the line $\delta(p)$:
\[
\pnt xy1 \sim \lin x1{-y} \quad (x \ne 0,\ y \notin S), \qquad \pnt x10 \sim \lin x0{-1} = \lin 10{-1/x} \quad (x \neq 0).
\]
\end{enumerate}
Let $G = G(q;a,c)$ be the incidence graph of $\cS$ (steps \ref{c:plane}--\ref{c:links}) and $R = R(q;a,c)$ the graph with the filler~\ref{c:fill} added. Both are bipartite with $2(q^2+q+1+5) = 2(q^2+q+6)$ vertices.

\paragraph{The rotation as a map.} By~\ref{f:lam}, in $\PG(2,q)$ the incidences between the points of $\ell$ and the lines through $A$ form the perfect matching $X_t \leftrightarrow \lam_t$, $P \leftrightarrow \lam_\infty$. Let $\pi\colon \F \to \F$ be $\sigma$ on $T$ and the identity on $\F \setminus T$. Step~\ref{c:plane} replaces this matching by $X_t \leftrightarrow \lam_{\pi(t)}$ ($t \in \F \setminus \{a\}$), $P \leftrightarrow \lam_\infty$, and leaves $X_a$ and $\lam_a$ unmatched. All other incidences of $\PG(2,q)$ are kept. (In $\cA$, $\pi$ is the identity and again only $X_a \leftrightarrow \lam_a$ is deleted.)

\paragraph{Changes relative to $\cA$.} Take $\cA$ with the same $a$ and choose $c \in \F \setminus \{0,1,a\}$. Then $\cS$ is obtained from $\cA$ by exactly these changes:
\begin{enumerate}[label=\textup{($\Delta$\arabic*)},leftmargin=3em]
\item add the point $H_c$, on the $q$ lines $\lin w1{-c}$ ($w \in \F$), on $M$ and on $K_1$; add the line $K_c$, through the $q$ points $\pnt xc1$ ($x \in \F$), through $Q$ and through $H_0$;
\item remove $X_0 \in \lam_0$, $X_1 \in \lam_1$, $X_c \in \lam_c$, and add $X_0 \in \lam_c$, $X_c \in \lam_1$, $X_1 \in \lam_0$;
\item remove $X_c \in M$ and $Q \in \lam_c$;
\item remove $H_0 \in K_1$ (it is replaced by $H_0 \in K_c$ and $H_c \in K_1$ from the first item; $H_1 \in K_0$ is kept).
\end{enumerate}
The core gains $2(q+2)$ incidences at $H_c$ and $K_c$ (they are not incident with each other) and three from the rotation, and loses six; that is, $2q+1$ more incidences than $\cA$. If both structures use the filler $p \sim \delta(p)$, the fillers differ only in that the points $\pnt xc1$ and the lines $\lin x1{-c}$, $x \neq 0$, are no longer matched (they now meet $K_c$ and $H_c$).

\subsection{Degrees}

\begin{table}[ht]
\centering\small
\renewcommand{\arraystretch}{1.15}
\begin{tabular}{@{}lll@{}}
\toprule
vertex & neighbours in $G$ & degree\\
\midrule
$A$ & $\lam_t$ ($t \in \F$), $\lam_\infty$, $K_a$ & $q+2$\\
$P$ & $\ell$, $\lam_\infty$, $\lin 10w$ ($w \in \F^*$), $M$ & $q+2$\\
$X_t$, $t \notin S$ & $\ell$, $\lin w1{-t}$ ($w \in \F$), $M$ & $q+2$\\
$X_a$ & $\ell$, $\lin w1{-a}$ ($w \in \F^*$), $M$, $K_a$ & $q+2$\\
$X_t$, $t \in T$ & $\ell$, $\lin w1{-t}$ ($w \in \F^*$), $\lam_{\sigma(t)}$, $K_t$ & $q+2$\\
$\pnt xy1$, $x \ne 0$ & $\lin u1{-(ux+y)}$ ($u \in \F$), $\lin 10{-x}$, and $K_y$ if $y \in S$ & $q+2$ if $y \in S$, else $q+1$\\
$\pnt x10$, $x \ne 0$ & $\lam_\infty$, $\lin 1{-x}w$ ($w \in \F$) & $q+1$\\
$Q$ & $\lam_\infty$, $\lam_t$ ($t \in \F\setminus T$), $K_s$ ($s \in S$) & $q+2$\\
$H_a$ & $\lin w1{-a}$ ($w \in \F$), $M$, $\ell$ & $q+2$\\
$H_t$, $t \in T$ & $\lin w1{-t}$ ($w \in \F$), $M$, $K_{\sigma(t)}$ & $q+2$\\
\midrule
$\ell$ & $P$, $X_t$ ($t \in \F$), $H_a$ & $q+2$\\
$\lam_\infty$ & $A$, $\pnt x10$ ($x \in \F$), $Q$ & $q+2$\\
$\lam_t$, $t \notin S$ & $A$, $\pnt xt1$ ($x \in \F$), $Q$ & $q+2$\\
$\lam_a$ & $A$, $\pnt xa1$ ($x \in \F^*$), $Q$, $H_a$ & $q+2$\\
$\lam_t$, $t \in T$ & $A$, $\pnt xt1$ ($x \in \F^*$), $X_{\sigma^{-1}(t)}$, $H_t$ & $q+2$\\
$\lin u1w$, $u \ne 0$ & $\pnt x{-ux-w}1$ ($x \in \F$), $\pnt{-1/u}10$, and $H_{-w}$ if $-w \in S$ & $q+2$ if $-w \in S$, else $q+1$\\
$\lin 10w$, $w \ne 0$ & $\pnt{-w}y1$ ($y \in \F$), $P$ & $q+1$\\
$M$ & $P$, $X_t$ ($t \in \F \setminus T$), $H_s$ ($s \in S$) & $q+2$\\
$K_a$ & $\pnt xa1$ ($x \in \F$), $Q$, $A$ & $q+2$\\
$K_t$, $t \in T$ & $\pnt xt1$ ($x \in \F$), $Q$, $H_{\sigma^{-1}(t)}$ & $q+2$\\
\bottomrule
\end{tabular}
\caption{The neighbourhood of every vertex of $G = G(q;a,c)$. Here $T = \{0,1,c\}$, $S = T \cup \{a\}$ and $\sigma = (0\,c\,1)$, so $\sigma^{-1}(0) = 1$, $\sigma^{-1}(1) = c$, $\sigma^{-1}(c) = 0$. The filler~\ref{c:fill} adds one more neighbour to exactly the vertices of degree $q+1$.}
\label{tab:nbrs}
\end{table}

\begin{lemma}[Degrees]\label{lem:coords-deg}
Table~\ref{tab:nbrs} lists the neighbourhood of every vertex of $G$. In particular $G$ has maximum degree $q+2$, and its vertices of degree $q+1$ are exactly the points $\pnt xy1$ \textup{(}$x \ne 0$, $y \notin S$\textup{)} and $\pnt x10$ \textup{(}$x \ne 0$\textup{)}, and the lines $\lin u1w$ \textup{(}$u \ne 0$, $-w \notin S$\textup{)} and $\lin 10w$ \textup{(}$w \ne 0$\textup{)}; there are $(q-1)(q-3)$ of each. The graph $R$ is $(q+2)$-regular.
\end{lemma}

\begin{proof}
Each row is \ref{f:lam} or~\ref{f:X} (the plane neighbours) corrected by \ref{c:plane}--\ref{c:links}. We check the rows one class at a time; in each case the listed neighbours are distinct, so the degree is their number.

\emph{Points.} $A$: the $q+1$ lines through $A$ are $\lam_t$ and $\lam_\infty$; no incidence of $A$ is removed; $A$ gains $K_a$ by~\ref{c:links} and nothing else ($A$ is not of the form $\pnt xs1$ in~\ref{c:hatline}). $P$: \ref{f:X}, no removal, gains $M$. $X_t$, $t \notin S$: \ref{f:X}, no removal, gains $M$ (as $t \notin T$), and no $K_s$ since $\pnt 0t1 = \pnt xs1$ forces $s = t$. $X_a$: loses $\lam_a$ (the case $w = 0$), gains $M$ ($a \notin T$) and $K_a$ (\ref{c:hatline} with $x=0$). $X_t$, $t \in T$: loses $\lam_t$, gains $\lam_{\sigma(t)}$ by~\ref{c:plane} and $K_t$; it is not on $M$. A point $\pnt xy1$ with $x \ne 0$: its plane lines are $\lin u1{-(ux+y)}$ and $\lin 10{-x}$ (solve $ux + v y + w = 0$ with $v = 1$, resp.\ $v = 0$); it is not on $\ell$, so it loses nothing, and it gains $K_s$ exactly when $y = s \in S$. A point $\pnt x10$ with $x \ne 0$: its lines are $\lam_\infty$ and $\lin 1{-x}w$, and it gains nothing. $Q$: $1 + (q-3) + 4 = q+2$ lines, by~\ref{c:Q}. $H_s$: $q$ lines by~\ref{c:hat}, $M$ by~\ref{c:M}, and exactly one line from~\ref{c:links}: $\ell$ if $s = a$, $K_{\sigma(s)}$ if $s \in T$.

\emph{Lines.} $\ell$: no removal (the removed incidences are on lines $\lam_s$), gains $H_a$. $\lam_\infty$: no removal, gains $Q$. $\lam_t$, $t \notin S$: no change except $Q$. $\lam_a$: loses $X_a$, gains $Q$ ($a \notin T$) and $H_a$ (\ref{c:hat} with $w = 0$). $\lam_t$, $t \in T$: loses $X_t$, gains $X_{\sigma^{-1}(t)}$ by~\ref{c:plane} and $H_t$; it does not contain $Q$. A line $\lin u1w$ with $u \neq 0$: its points are $\pnt x{-ux-w}1$ and $\pnt{-1/u}10$; it is not a line through $A$, so it loses nothing; it gains $H_s$ exactly when $\lin u1w = \lin{u}1{-s}$, i.e.\ $s = -w \in S$. A line $\lin 10w$, $w \ne 0$: its points are $\pnt{-w}y1$ and $P$, and it gains nothing (it is not of the form $\lin{w'}1{-s}$). $M$: $1 + (q-3) + 4 = q+2$ points. $K_s$: $q$ points by~\ref{c:hatline}, $Q$, and exactly one point from~\ref{c:links}: $A$ if $s = a$, $H_{\sigma^{-1}(s)}$ if $s \in T$.

\emph{Filler.} The points of degree $q+1$ number $(q-1)(q-4) + (q-1) = (q-1)(q-3)$, and so do the lines. The map $\delta$ sends $\pnt xy1 \mapsto \lin x1{-y}$, a line $\lin u1w$ with $u = x \ne 0$ and $-w = y \notin S$, and $\pnt x10 \mapsto \lin x0{-1}$, a line $\lin 10{w}$ with $w = -1/x \ne 0$. It is injective, so it is a bijection between the points and the lines of degree $q+1$. Since these points are not on $\ell$, $p \notin \delta(p)$ (\S\ref{ssec:notation}), so~\ref{c:fill} adds new edges forming a perfect matching between the vertices of degree $q+1$.
\end{proof}

\subsection{Diameter}

\begin{lemma}\label{lem:coords-reduce}
Let $\mathcal{I}$ be a finite incidence structure in which every line has at least one point, any two points lie on a common line and any two lines meet in a common point. Then its incidence graph has diameter at most $3$.
\end{lemma}

\begin{proof}
Two points, or two lines, are at distance at most $2$. Let $p$ be a point and $L$ a line with $p \notin L$. Take a point $p' \in L$; then $p' \ne p$. Take a line $L'$ through $p$ and $p'$; then $L' \ne L$, and $p, L', p', L$ is a path.
\end{proof}

In the next two lemmas, \emph{incident} means incident in $\cS$. Recall that the only incidences of $\PG(2,q)$ missing from $\cS$ are $X_s \in \lam_s$, $s \in S$. So an incidence $u \in L$ of $\PG(2,q)$ holds in $\cS$ whenever $u \notin \ell$ or $L$ is not a line through $A$.

\begin{lemma}[Points]\label{lem:coords-points}
Any two points of $\cS$ lie on a common line of $\cS$.
\end{lemma}

\begin{proof}
Let $u \ne v$ be points.

\emph{(P1) Two points of $\PG(2,q)$.} Let $L$ be the line of $\PG(2,q)$ through $u$ and $v$. If both incidences hold in $\cS$, use $L$. Otherwise, say, $u = X_s$ and $L = \lam_s$ with $s \in S$, and by~\ref{f:lam} $v = A$ or $v = \pnt xs1$ with $x \ne 0$.
If $v = \pnt xs1$: $K_s$ contains $u = \pnt 0s1$ and $v$ by~\ref{c:hatline}.
If $v = A$ and $s = a$: $K_a$ contains $X_a$ (\ref{c:hatline}, $x=0$) and $A$ (\ref{c:links}).
If $v = A$ and $s \in T$: $\lam_{\sigma(s)}$ contains $X_s$ (\ref{c:plane}) and $A$.

\emph{(P2) $Q$ and a point $u$ of $\PG(2,q)$.}
If $u \in \{A, P\}$ or $u = \pnt x10$ ($x \ne 0$): $\lam_\infty$, by~\ref{f:lam} and~\ref{c:Q}.
If $u = X_t$, $t \notin S$: $\lam_t$ ($t \notin T$, and $X_t \in \lam_t$ is kept).
If $u = X_s$, $s \in S$: $K_s$ (\ref{c:hatline} with $x = 0$, \ref{c:Q}).
If $u = \pnt xy1$, $x \ne 0$: $u \in \lam_y$ by~\ref{f:lam}; use $\lam_y$ if $y \notin T$, and $K_y$ if $y \in T$.

\emph{(P3) $H_s$ and a point $u$ of $\PG(2,q)$.}
If $u \notin \ell$ (this includes $A$), write $u = \pnt xyz$ with $x \neq 0$: by~\ref{f:join} $u$ lies on $\lin{(sz-y)/x}1{-s}$, which contains $H_s$ by~\ref{c:hat}.
If $u = X_s$: the line $\lin 11{-s}$ contains $X_s$, is not $\lam_s$, and contains $H_s$.
If $u = P$ or $u = X_t$ with $t \notin T$, $t \neq s$: $M$ (\ref{c:M}).
If $u = X_t$ with $t \in T$, $t \ne s$, and $s = a$: $\ell$ (\ref{c:links}).
If $u = X_t$ with $t \in T$, $t \neq s$, and $s \in T$: since $\sigma$ is a $3$-cycle on $T$, either $t = \sigma(s)$ or $t = \sigma^{-1}(s)$. If $t = \sigma(s)$, use $K_{\sigma(s)}$, which contains $X_{\sigma(s)}$ (\ref{c:hatline}) and $H_s$ (\ref{c:links}). If $t = \sigma^{-1}(s)$, use $\lam_s = \lin 01{-s}$, which contains $X_{\sigma^{-1}(s)}$ (\ref{c:plane}) and $H_s$ (\ref{c:hat} with $w = 0$).

\emph{(P4) Two new points.} $H_s$ and $H_{s'}$: $M$. $Q$ and $H_a$: $\lam_a$ ($Q$ by~\ref{c:Q} as $a \notin T$; $H_a$ by~\ref{c:hat} with $w=0$). $Q$ and $H_t$, $t \in T$: $K_{\sigma(t)}$ (\ref{c:Q}, \ref{c:links}).
\end{proof}

\begin{lemma}[Lines]\label{lem:coords-lines}
Any two lines of $\cS$ meet in a common point of $\cS$.
\end{lemma}

\begin{proof}
Let $L \ne N$ be lines.

\emph{(L1) Two lines of $\PG(2,q)$.} Let $p$ be their common point in $\PG(2,q)$. If both incidences hold in $\cS$, use $p$. Otherwise, say, $p = X_s$ and $L = \lam_s$ with $s \in S$, and by~\ref{f:X} $N = \ell$ or $N = \lin w1{-s}$ with $w \ne 0$.
If $N = \lin w1{-s}$: $H_s$ lies on $N$ and on $\lam_s = \lin 01{-s}$ by~\ref{c:hat}.
If $N = \ell$ and $s = a$: $H_a$ lies on $\lam_a$ (\ref{c:hat}, $w = 0$) and on $\ell$ (\ref{c:links}).
If $N = \ell$ and $s \in T$: $X_{\sigma^{-1}(s)}$ lies on $\ell$ and on $\lam_s$ (\ref{c:plane}).

\emph{(L2) $M$ and a line $L$ of $\PG(2,q)$.}
If $L$ passes through $P$, i.e.\ $L \in \{\ell, \lam_\infty\}$ or $L = \lin 10w$ ($w \ne 0$) by~\ref{f:X}: $P$.
If $L = \lam_t$, $t \notin S$: $X_t$ ($t \notin T$, and $X_t \in \lam_t$ is kept).
If $L = \lam_s$, $s \in S$: $H_s$ (\ref{c:hat} with $w = 0$, \ref{c:M}).
If $L = \lin u1w$ with $u \ne 0$: by~\ref{f:meet} $L$ meets $\ell$ in $X_r$, $r = -w$, and this incidence is kept. If $r \notin T$, use $X_r \in M$. If $r \in T$, then $L = \lin u1{-r}$ contains $H_r$ by~\ref{c:hat}, and $H_r \in M$.

\emph{(L3) $K_s$ and a line $L$ of $\PG(2,q)$.}
If $L$ does not pass through $A$ (that is, $L = \ell$, or $L = \lin u v w$ with $u \neq 0$): $L$ meets $\lam_s$ in a point $\pnt xs1$ (by~\ref{f:lam} for $L = \ell$, where the point is $X_s$; by~\ref{f:meet} otherwise), the incidence with $L$ is kept, and the point is on $K_s$ by~\ref{c:hatline}.
If $L = \lam_s$: the point $\pnt 1s1$ (on $\lam_s$, not on $\ell$) and $K_s$.
If $L = \lam_\infty$ or $L = \lam_t$ with $t \notin T$: $Q$ (\ref{c:Q}).
If $L = \lam_t$ with $t \in T$, $t \ne s$, and $s = a$: $A$ (\ref{c:links}).
If $L = \lam_t$ with $t \in T$, $t \ne s$, and $s \in T$: if $t = \sigma(s)$, the point $X_s$ lies on $\lam_{\sigma(s)}$ (\ref{c:plane}) and on $K_s$ (\ref{c:hatline}); if $t = \sigma^{-1}(s)$, the point $H_t$ lies on $\lam_t$ (\ref{c:hat}, $w = 0$) and on $K_{\sigma(t)} = K_s$ (\ref{c:links}).

\emph{(L4) Two new lines.} $K_s$ and $K_{s'}$: $Q$. $M$ and $K_a$: $X_a$ (\ref{c:M} as $a \notin T$; \ref{c:hatline} with $x = 0$). $M$ and $K_t$, $t \in T$: $H_{\sigma^{-1}(t)}$ (\ref{c:M}, \ref{c:links}).
\end{proof}

\begin{theorem}\label{thm:coords}
Let $q \ge 4$ be a prime power and $a, c$ distinct elements of $\mathrm{GF}(q) \setminus \{0,1\}$. The graph $G(q;a,c)$ is bipartite, has $2(q^2+q+6)$ vertices, maximum degree $q+2$ and diameter $3$. The graph $R(q;a,c)$ has the same properties and is $(q+2)$-regular.
\end{theorem}

\begin{proof}
The order and the degrees are Lemma~\ref{lem:coords-deg}. Every line of $\cS$ has at least one point. By Lemmas~\ref{lem:coords-points} and~\ref{lem:coords-lines} and Lemma~\ref{lem:coords-reduce}, $G$ has diameter at most $3$; adding the edges~\ref{c:fill} cannot increase it. The diameter is at least $3$ because the graphs are bipartite and not complete bipartite (for instance $A \notin \ell$).
\end{proof}

\begin{remark}[What the proofs use]
Only that $\F$ is a field and that $0, 1, a, c$ are distinct: the divisions in \ref{f:join}, \ref{f:meet} and Table~\ref{tab:nbrs} are by $x \ne 0$ or $u \ne 0$, and the filler uses $x^2 \neq 0$ for $x \neq 0$. No assumption on the characteristic is needed. The values $0, 1, c$ enter only through $T$ and the $3$-cycle $\sigma$, so the same proofs hold verbatim for any three distinct elements $t_0, t_1, t_2 \in \F \setminus \{a\}$ and $\sigma = (t_0\,t_2\,t_1)$. The point $P$ plays the same role as every $X_t$ with $t \notin S$; it is not a parameter.
\end{remark}

\begin{remark}[Duality]\label{rem:duality}
Write $\cS_\sigma$ for the structure built with the $3$-cycle $\sigma$ on $T$. The correlation $\delta$ maps the points of $\cS_\sigma$ to the lines of $\cS_{\sigma^{-1}}$ and vice versa, preserving incidence, if we extend it by $\delta(H_s) = K_s$, $\delta(K_s) = H_s$, $\delta(Q) = M$, $\delta(M) = Q$. For instance, $\delta$ maps $X_t \in \lam_{\sigma(t)}$ to $X_{\sigma(t)} \in \lam_t$, which is the rotation of $\cS_{\sigma^{-1}}$, and maps $\pnt xs1 \in K_s$ to $H_s \in \lin x1{-s}$. So Lemma~\ref{lem:coords-lines} also follows from Lemma~\ref{lem:coords-points} applied to $\cS_{\sigma^{-1}}$; we have written it out so that each case can be checked directly. For $\cA$, $\delta$ is a correlation of $\cA$ itself.
\end{remark}

\begin{remark}[Computer check]
The program \texttt{coords.py} builds $\cA$ and $\cS$ only from this section, in three ways (from the steps \ref{c:plane}--\ref{c:links}, from the point rows of Table~\ref{tab:nbrs}, and from its line rows), checks that the three incidence sets coincide, and checks the degrees, both covering properties and the exact diameter of $G$ and $R$, for every prime power $q \le 13$ and every admissible $(a,c)$: $2, 6, 20, 30, 42, 72, 110$ choices for $q = 4, 5, 7, 8, 9, 11, 13$. All pass. The cores $G$ are isomorphic (by \texttt{nauty}) to those of the earlier coordinate-free implementation for $q = 9, 11, 13$, and the graphs $R$ for $q = 9, 11, 13$ pass three independent diameter checkers. It also builds $\cA$ for every prime power $q \le 13$ and every $a$, confirms the degree count of \S\ref{ssec:al} (with $i \in \F \setminus \{0,1,a\}$ the matching of~\cite{AL} leaves four vertices of degree $q+1$; with $i \in \F^*$ the graph is regular), and finds the graphs isomorphic to an independent implementation of~\cite{AL}. It checks the duality of Remark~\ref{rem:duality} directly. Removing the rotation, or the incidence $H_1 \in K_0$, gives diameter $4$ in every case tested ($q = 5, 9, 13$).
\end{remark}
